\documentclass[11pt,fontset=fandol]{ctexart}
\usepackage[a4paper,margin=23mm]{geometry}
\usepackage{amsmath,amssymb,mathtools,bm,mathrsfs,esint}
\usepackage{graphicx,longtable,booktabs,array,calc}
\usepackage[normalem]{ulem}
\usepackage{xcolor,hyperref}
\hypersetup{hidelinks}
\setcounter{secnumdepth}{-1}
\newcounter{none}
\providecommand{\tightlist}{\setlength{\itemsep}{0pt}\setlength{\parskip}{0pt}}
\providecommand{\passthrough}[1]{#1}
\usepackage{amssymb,mathrsfs,esint,bm,newunicodechar}
\allowdisplaybreaks
\setlength{\emergencystretch}{3em}
\providecommand{\llangle}{\langle\!\langle}
\providecommand{\rrangle}{\rangle\!\rangle}
\newunicodechar{∞}{\ensuremath{\infty}}
\newunicodechar{∈}{\ensuremath{\in}}
\newunicodechar{∑}{\ensuremath{\sum}}
\newunicodechar{∫}{\ensuremath{\int}}
\newunicodechar{≤}{\ensuremath{\leq}}
\newunicodechar{≥}{\ensuremath{\geq}}
\newunicodechar{≠}{\ensuremath{\neq}}
\newunicodechar{→}{\ensuremath{\to}}
\newunicodechar{⇒}{\ensuremath{\Rightarrow}}
\newunicodechar{α}{\ensuremath{\alpha}}
\newunicodechar{β}{\ensuremath{\beta}}
\newunicodechar{γ}{\ensuremath{\gamma}}
\newunicodechar{λ}{\ensuremath{\lambda}}
\newunicodechar{φ}{\ensuremath{\varphi}}
\newunicodechar{ε}{\ensuremath{\varepsilon}}
\newunicodechar{σ}{\ensuremath{\sigma}}
\newunicodechar{θ}{\ensuremath{\theta}}
\newunicodechar{δ}{\ensuremath{\delta}}
\newunicodechar{η}{\ensuremath{\eta}}
\newunicodechar{π}{\ensuremath{\pi}}
\newunicodechar{ξ}{\ensuremath{\xi}}
\newunicodechar{Σ}{\ensuremath{\Sigma}}
\newunicodechar{Δ}{\ensuremath{\Delta}}
\newunicodechar{∂}{\ensuremath{\partial}}
\newunicodechar{∀}{\ensuremath{\forall}}
\newunicodechar{∃}{\ensuremath{\exists}}
\newunicodechar{√}{\ensuremath{\surd}}
\newunicodechar{−}{\ensuremath{-}}
\newunicodechar{μ}{\ensuremath{\mu}}
\newunicodechar{ρ}{\ensuremath{\rho}}
\newunicodechar{τ}{\ensuremath{\tau}}
\newunicodechar{ω}{\ensuremath{\omega}}
\newunicodechar{Ω}{\ensuremath{\Omega}}
\newunicodechar{χ}{\ensuremath{\chi}}
\newunicodechar{ψ}{\ensuremath{\psi}}
\newunicodechar{Γ}{\ensuremath{\Gamma}}
\newunicodechar{Φ}{\ensuremath{\Phi}}
\newunicodechar{Λ}{\ensuremath{\Lambda}}
\newunicodechar{ℝ}{\ensuremath{\mathbb{R}}}
\newunicodechar{ℤ}{\ensuremath{\mathbb{Z}}}
\newunicodechar{ℕ}{\ensuremath{\mathbb{N}}}
\newunicodechar{△}{\ensuremath{\triangle}}
\newunicodechar{⊗}{\ensuremath{\otimes}}
\begin{document}
\section{2018年杭州电子科技大学601数学分析考研真题参考答案}\label{ux5e74ux676dux5ddeux7535ux5b50ux79d1ux6280ux5927ux5b66601ux6570ux5b66ux5206ux6790ux8003ux7814ux771fux9898ux53c2ux8003ux7b54ux6848}

来源：\texttt{HDU/601数学分析/2018/2018年杭州电子科技大学601数学分析考研真题参考答案.pdf}。原
PDF 页序：1, 2, 3, 4, 5, 6, 7, 8, 9。

\section{2018年数学分析打印参考答案：OCR核定整理}\label{ux5e74ux6570ux5b66ux5206ux6790ux6253ux5370ux53c2ux8003ux7b54ux6848ocrux6838ux5b9aux6574ux7406}

\begin{quote}
以下以原识别稿为基础。标注``【OCR公式整理】''的段落由助手依据同题题面、原解题思路与可核算公式恢复，不作为原作者逐字转录；识读分歧另行注明。
\end{quote}

\subsection{杭州电子科技大学2018年数学分析真题解析}\label{ux676dux5ddeux7535ux5b50ux79d1ux6280ux5927ux5b662018ux5e74ux6570ux5b66ux5206ux6790ux771fux9898ux89e3ux6790}

\subsection{一.填空题(本大题共8小题，每小题5分，本大题共40分)}\label{ux4e00.ux586bux7a7aux9898ux672cux5927ux9898ux51718ux5c0fux9898ux6bcfux5c0fux98985ux5206ux672cux5927ux9898ux517140ux5206}

\begin{enumerate}
\def\labelenumi{\arabic{enumi}.}
\tightlist
\item
  \(\lim _ { n  \infty } ( 2 ^ { n } + 3 ^ { n } + 4 ^ { n } ) ^ { \frac { 1 } { n } } = (\)
\end{enumerate}

【宝哥解答】

【OCR公式整理】设 \(a_1,\ldots,a_k\) 为非负常数，\(M=\max_i a_i\)。若
\(M=0\)，结论显然；若 \(M>0\)，则 \[
M\le\sqrt[n]{a_1^n+\cdots+a_k^n}\le k^{1/n}M\longrightarrow M.
\] 故本题极限为 \(4\)。根号次数和夹逼末式依据同一不等式恢复。

【题面符号待核对】独立真题OCR将条件识为 \(f(1)=3\)，当前参考稿明确写作
\(f'(1)=3\)。一般极限等于 \(4f'(1)\)；数值 \(12\)
仅对应后一条件，若只有函数值则不能确定。

2.设函数f（x）在x=1处可导， \(x = 1\) \(f ^ { \prime } ( 1 ) = 3\) ，则
\(\lim _ { x \to 0 } { \frac { f ( 1 + x ) - f ( 1 - 3 x ) } { x } } = (\)

\[
\begin{array} { l } { \displaystyle \lim _ { x \to 0 } \frac { f ( 1 + x ) - f ( 1 - 3 x ) } { x } = \lim _ { x \to 0 } \left[ \frac { f ( 1 + x ) - f ( 1 ) } { x } + 3 \frac { f ( 1 - 3 x ) - f ( 1 ) } { - 3 x } \right] } \\ { = f ^ { \prime } ( 1 ) + 3 f ^ { \prime } ( 1 ) = 4 f ^ { \prime } ( 1 ) = 4 \times 3 = 1 2 } \end{array}
\]

3.已知 \(e ^ { x + y } = \sin ( x y )\) 确定了隐函数y=y(x)，则
\(\frac { d y } { d x } =\) 研习

\subsection{【宝哥解答】}\label{ux5b9dux54e5ux89e3ux7b54}

令 \(F ( x , y ) = e ^ { x + y } - \sin ( x y )\)
，由隐函数存在定理的结论，

\[
\frac { d y } { d x } = - \frac { F _ { x } } { F _ { y } } = - \frac { e ^ { x + y } - y \cos ( x y ) } { e ^ { x + y } - x \cos ( x y ) }
\]

4.设 \(F ( x ) = \int _ { 0 } ^ { 2 x } x e ^ { t ^ { 2 } } d t\) ，则
\(F ^ { \prime } ( x ) = (\)

\subsection{【宝哥解答】}\label{ux5b9dux54e5ux89e3ux7b54-1}

\[
F ( x ) = \int _ { 0 } ^ { 2 x } x e ^ { t ^ { 2 } } d t = x \int _ { 0 } ^ { 2 x } e ^ { t ^ { 2 } } d t \ ,
\]

\[
F ^ { \prime } ( x ) = \int _ { 0 } ^ { 2 x } e ^ { t ^ { 2 } } d t + x e ^ { ( 2 x ) ^ { 2 } } \cdot 2 = \int _ { 0 } ^ { 2 x } e ^ { t ^ { 2 } } d t + 2 x e ^ { 4 x ^ { 2 } }
\]

5.设 \(u = x ^ { 2 } + y - z ^ { 2 }\)
，则u在(1,-1,1)处的方向导数的最大值为(

\subsection{【宝哥解答】}\label{ux5b9dux54e5ux89e3ux7b54-2}

\[
u = x ^ { 2 } + y - z ^ { 2 } , u _ { x } = 2 x , u _ { y } = 1 , u _ { z } = - 2 z ,
\]

都是连续的，故u在(x,y,z)处的方向导数的最大值就是该点梯度的模，为\(\sqrt { u _ { x } ^ { 2 } + u _ { y } ^ { 2 } + u _ { z } ^ { 2 } } = \sqrt { 4 x ^ { 2 } + 1 + 4 z ^ { 2 } }\)
，故u在(1,-1,1)处的方向导数的最大值为\(\sqrt { 4 \times 1 ^ { 2 } + 1 + 4 \times 1 ^ { 2 } } = 3\)

\begin{enumerate}
\def\labelenumi{\arabic{enumi}.}
\setcounter{enumi}{5}
\tightlist
\item
  已知 \((x+ay)\,dx+(2x+y)\,dy\) 为某函数的全微分，求 \(a\)。
\end{enumerate}

【OCR公式整理】令 \(P=x+ay\)、\(Q=2x+y\)。全微分条件给出
\(P_y=Q_x\)，所以 \(a=2\)。

7.设函数f（x，y）连续，则二次积分
\(\int _ { - 2 } ^ { 0 } d x \int _ { 2 x + 4 } ^ { 4 - x ^ { 2 } } f ( x , y ) d y\)
可交换积分次序为

【宝哥解答】

作出积分区域如下图所示：

于是，

\[
\int _ { - 2 } ^ { 0 } d x \int _ { 2 x + 4 } ^ { 4 - x ^ { 2 } } f ( x , y ) d y = \int _ { 0 } ^ { 4 } d y \int _ { - \sqrt { 4 - y } } ^ { \frac { y - 4 } { 2 } } f ( x , y ) d x
\]

\begin{enumerate}
\def\labelenumi{\arabic{enumi}.}
\setcounter{enumi}{7}
\tightlist
\item
  设 \(\Sigma:x^2+y^2+z^2=1\)，求曲面积分 \[
  \iint_\Sigma(2-x^2-y^2-z^2)\,dS.
  \] 【OCR公式整理】球面上被积函数恒为 \(1\)，单位球面面积为
  \(4\pi\)，故积分为 \(4\pi\)。
\end{enumerate}

\subsection{二.计算题(本大题共5小题，每小题6分，本大题共30分)}\label{ux4e8c.ux8ba1ux7b97ux9898ux672cux5927ux9898ux51715ux5c0fux9898ux6bcfux5c0fux98986ux5206ux672cux5927ux9898ux517130ux5206}

\begin{enumerate}
\def\labelenumi{\arabic{enumi}.}
\tightlist
\item
  求
  \(\displaystyle\lim_{x\to\infty}\frac{x^3(1-\cos(1/x^2))}{\sqrt{x^2+1}-x}\)。
\end{enumerate}

【OCR公式整理】按原参考解答分别讨论正、负无穷，避免混淆方向。由于
\(1-\cos(1/x^2)\sim1/(2x^4)\)，当 \(x\to+\infty\) 时， \[
\frac{x^3(1-\cos(1/x^2))}{\sqrt{x^2+1}-x}
\sim\frac{\sqrt{x^2+1}+x}{2x}\longrightarrow1.
\] 当 \(x\to-\infty\) 时，分母 \(\sim-2x\)，所以原式
\(\sim-1/(4x^2)\to0\)。 因此，若题面的 \(\infty\) 指
\(|x|\to\infty\)，极限不存在；若指 \(+\infty\)，答案为
\(1\)。原卷极限方向符号需核对。

\begin{enumerate}
\def\labelenumi{\arabic{enumi}.}
\setcounter{enumi}{1}
\tightlist
\item
  设 \(u=u(x,y)\)、\(v=v(x,y)\) 由隐函数组确定，求 \(u_x\)。
\end{enumerate}

【题面互证及OCR整理】同年真题识别稿与参考解答的求导步骤均支持 \[
\begin{cases}\sin(u+v)=xy,\\e^{uv}=1-x-y.\end{cases}
\] 参考答案题面中另有 \(1-x^2y\)
的识别，不能与其后原求导过程相容；此处按前述互证版本计算，原稿该处仍需核对。对
\(x\) 求偏导得 \[
\cos(u+v)(u_x+v_x)=y,\qquad
e^{uv}(vu_x+uv_x)=-1.
\] 在隐函数条件 \((u-v)e^{uv}\cos(u+v)\ne0\) 下， \[
u_x=\frac{yu e^{uv}+\cos(u+v)}{(u-v)e^{uv}\cos(u+v)}.
\] 同理，全微分为 \[
du=\frac{[yu e^{uv}+\cos(u+v)]\,dx+[xu e^{uv}+\cos(u+v)]\,dy}
{(u-v)e^{uv}\cos(u+v)}.
\]

\begin{enumerate}
\def\labelenumi{\arabic{enumi}.}
\setcounter{enumi}{2}
\tightlist
\item
  已知连续函数 \(f\) 满足 \(\int_0^x tf(x-t)\,dt=1-\cos x\)，求
  \(\int_0^{\pi/2}f(x)\,dx\)。
\end{enumerate}

【OCR公式整理】令 \(s=x-t\)，则 \[
\int_0^x tf(x-t)\,dt
=x\int_0^x f(s)\,ds-\int_0^x sf(s)\,ds.
\] 对 \(x\) 求导得 \(\int_0^x f(s)\,ds=\sin x\)，故所求积分为 \(1\)。

\begin{enumerate}
\def\labelenumi{\arabic{enumi}.}
\setcounter{enumi}{3}
\tightlist
\item
  求 \(\displaystyle\sum_{n=0}^\infty\frac{x^n}{n+1}\) 的和函数。
\end{enumerate}

【OCR公式整理】收敛域为 \([-1,1)\)。对 \(|x|<1\)，先对几何级数逐项积分：
\[
\sum_{n=0}^\infty\frac{x^{n+1}}{n+1}
=\int_0^x\frac{dt}{1-t}=-\ln(1-x).
\] 故 \[
S(x)=\begin{cases}
-\ln(1-x)/x,&-1\le x<1,\ x\ne0,\\
1,&x=0.
\end{cases}
\] 其中 \(x=-1\) 的值由交错级数及 Abel 定理确定，为
\(\ln2\)。原识别中分母误带阶乘和遗漏的最终和函数据此恢复。

\begin{enumerate}
\def\labelenumi{\arabic{enumi}.}
\setcounter{enumi}{4}
\tightlist
\item
  计算 \(\iiint_\Omega z^2\,dV\)，其中 \(\Omega\) 由 \(z=x^2+y^2\) 与
  \(z=\sqrt{x^2+y^2}\) 围成。
\end{enumerate}

【OCR公式整理】水平截面为
\(z^2\le x^2+y^2\le z\)，\(0\le z\le1\)，面积为 \(\pi(z-z^2)\)。因此 \[
\iiint_\Omega z^2\,dV=\pi\int_0^1 z^2(z-z^2)\,dz
=\pi\left(\frac14-\frac15\right)=\frac{\pi}{20}.
\]

\subsection{三.综合题（本大题共3题，每题10分，本大题共30分）}\label{ux4e09.ux7efcux5408ux9898ux672cux5927ux9898ux51713ux9898ux6bcfux989810ux5206ux672cux5927ux9898ux517130ux5206}

\begin{enumerate}
\def\labelenumi{\arabic{enumi}.}
\tightlist
\item
  设 \(f(0)=0\)、\(f''(0)\) 存在， \[
  g(x)=\begin{cases}f(x)/x,&x\ne0,\\f'(0),&x=0.\end{cases}
  \] 求 \(g'(0)\)。
\end{enumerate}

【OCR公式整理】分段函数在零点的定义由同年真题识别稿恢复。由二阶 Taylor
展开， \[
f(x)=f'(0)x+\frac12f''(0)x^2+o(x^2),
\qquad g'(0)=\lim_{x\to0}\frac{f(x)-f'(0)x}{x^2}
=\frac12f''(0).
\]

\begin{enumerate}
\def\labelenumi{\arabic{enumi}.}
\setcounter{enumi}{1}
\tightlist
\item
  讨论 \(\sum_{n=1}^\infty(-1)^n(\sin^2x+n)/n^2\)
  在全实轴上的一致收敛性和绝对收敛性。
\end{enumerate}

【OCR公式整理】分解为 \[
\sum_{n=1}^\infty\frac{(-1)^n\sin^2x}{n^2}
+\sum_{n=1}^\infty\frac{(-1)^n}{n}.
\] 第一个级数由 \(1/n^2\)
控制，在全实轴上一致收敛；第二个是收敛的常数级数，故原级数一致收敛。对任意
\(x\)，通项的绝对值不小于 \(1/n\)，故处处不绝对收敛。

\begin{enumerate}
\def\labelenumi{\arabic{enumi}.}
\setcounter{enumi}{2}
\tightlist
\item
  计算曲面积分 \[
  \iint_\Sigma x^2\,dy\,dz+y^2\,dz\,dx+z(1-2x-2y)\,dx\,dy,
  \] 其中 \(\Sigma\) 为上半单位球面 \(x^2+y^2+z^2=1\)、\(z\ge0\)
  的上侧。
\end{enumerate}

【宝哥解答】

【OCR公式整理】补上 \(z=0\)、\(x^2+y^2\le1\) 的下侧圆盘
\(\Sigma_1\)，与上半球面共同组成外向闭曲面。向量场为 \[
\mathbf F=(x^2,y^2,z(1-2x-2y)),\qquad \operatorname{div}\mathbf F=1.
\] 圆盘上第三分量为零，通量为零。由 Gauss 公式，原积分为 \[
\iint_\Sigma\mathbf F\cdot\mathbf n\,dS
=\iiint_\Omega1\,dV=\frac{2\pi}{3}.
\]

\subsection{四.证明题(本大题共5题，每题10分，本大题共50分)
四.证明题（本大题共5题，每题10分，本人起六}\label{ux56db.ux8bc1ux660eux9898ux672cux5927ux9898ux51715ux9898ux6bcfux989810ux5206ux672cux5927ux9898ux517150ux5206-ux56db.ux8bc1ux660eux9898ux672cux5927ux9898ux51715ux9898ux6bcfux989810ux5206ux672cux4ebaux8d77ux516d}

1.设函数f（x）在 \(( 0 , + \infty )\) 上满足 \(f ( x ) = f ( 3 x )\) 且
\(\lim _ { x \to + \infty } f ( x ) = 5\) .证明在 \(( 0 , + \infty )\)
上 \(f ( x ) \equiv 5 .\)

\subsection{【宝哥解答】}\label{ux5b9dux54e5ux89e3ux7b54-3}

函数f（x）在（0,+∞0）上满足 \(f ( x ) = f ( 3 x )\)
\(\lim _ { x \to + \infty } f ( x ) = 5\) 任取
\(x \in ( 0 , + \infty )\) 有
\(f ( x ) = f ( 3 x ) = f ( 3 ^ { 2 } x ) = \cdots = f ( 3 ^ { n } x )\)
\(\lim _ { n \to \infty } ( 3 ^ { n } x ) = + \infty\) 故
\(\lim _ { n \to \infty } f ( 3 ^ { n } x ) = 5\)
，故\(f ( x ) = \lim _ { n \to \infty } f ( 3 ^ { n } x ) = 5\)
由x的任意性，在 \(( 0 , + \infty ) [ f ( x ) \equiv 5\)

\begin{enumerate}
\def\labelenumi{\arabic{enumi}.}
\setcounter{enumi}{1}
\tightlist
\item
  设 \[
  f(x,y)=\begin{cases}
  \dfrac{x(x^2+y^2)}{x^2+y^4},&(x,y)\ne(0,0),\\
  0,&(x,y)=(0,0).
  \end{cases}
  \] 证明原点极限不存在，而两个偏导数存在。
\end{enumerate}

【OCR公式整理】沿 \(x=0\)，函数恒为零；沿 \(x=y^2\)， \[
f(y^2,y)=\frac{y^2(y^4+y^2)}{2y^4}=\frac{1+y^2}{2}\to\frac12,
\] 故二元极限不存在。由定义， \[
f_x(0,0)=\lim_{h\to0}\frac{f(h,0)}h=1,\qquad
f_y(0,0)=\lim_{h\to0}\frac{f(0,h)}h=0.
\]

3.设连续函数列 \(\{ f _ { n } ( x ) \}\) 在闭区间
\(\textstyle { \left[ a , b \right] }\) 上一致收敛于f(x).若
\(\boldsymbol { x } _ { n } \in \left[ a , b \right]\)
\(( n = 1 , 2 , \cdots )\) 且有
\(\lim _ { n \to \infty } x _ { n } = x _ { 0 }\) ，证明
\(\lim _ { n \to \infty } f _ { n } ( x _ { n } ) = f ( x _ { 0 } )\)

\subsection{【宝哥解答】}\label{ux5b9dux54e5ux89e3ux7b54-4}

连续函数列 \(\left\{ f _ { r } ( x ) \right\}\)
在闭区间{[}a,b{]}上一致收敛于f(x)，故f(x)也是连续的.\(\left\{ f _ { n } ( x ) \right\}\)
在{[}a,b{]}上一致收敛于f(x)，故对任意 \(\varepsilon > 0\) ，存在
\(N > 0\) ，使得对任意 \(n > N\) 和 \(x \in \left[ a , b \right]\)
，恒有 \(| f _ { n } ( x ) - f ( x ) | < \frac { 1 } { 2 } \varepsilon\)
.数列 \(\{ x _ { n } \} \subset [ a , b ]\)
\(\lim _ { n \to \infty } x _ { n } = x _ { 0 }\) 故
\(\left| f _ { n } ( x _ { n } ) - f ( x _ { n } ) \right| < { \frac { 1 } { 2 } } \varepsilon _ { \ast }\)
故

\[
| f _ { n } ( x _ { n } ) { - } f ( x _ { 0 } ) | \leq | f _ { n } ( x _ { n } ) - f ( x _ { n } ) | + | f ( x _ { n } ) - f ( x _ { 0 } ) | < { \frac { 1 } { 2 } } \varepsilon + | f ( x _ { n } ) - f ( x _ { 0 } ) |
\]

f（x）连续， \(\lim _ { n \to \infty } x _ { n } = x _ { 0 }\) ，故
\(\lim _ { n \to \infty } \left[ f ( x _ { n } ) - f ( x _ { 0 } ) \right] = 0\)
，故存在 \(N _ { 1 } > N\) ，使得对任意\(n > N _ { 1 }\) ，恒有
\(\left| f ( x _ { n } ) - f ( x _ { 0 } ) \right| < { \frac { 1 } { 2 } } \varepsilon\)
，故

\[
\left| f _ { n } ( x _ { n } ) - f ( x _ { 0 } ) \right| < { \frac { 1 } { 2 } } \varepsilon + { \frac { 1 } { 2 } } \varepsilon = \varepsilon
\]

故 \(\lim _ { n \to \infty } f _ { n } ( x _ { n } ) = f ( x _ { 0 } )\)

\begin{enumerate}
\def\labelenumi{\arabic{enumi}.}
\setcounter{enumi}{3}
\tightlist
\item
  设 \(f,\varphi\) 在 \((x_0,y_0)\)
  某邻域内可微，\(\varphi_y(x_0,y_0)\ne0\)，且该点为约束 \(\varphi=0\)
  下的极值点。证明 \(f_x(x_0,y_0)\ne0\) 时，\(f_y(x_0,y_0)\ne0\)。
\end{enumerate}

【OCR公式整理】原题第5题题干被并入本题末尾，现按题号恢复。由约束极值的一阶必要条件，存在
\(\lambda\)，使 \[
f_x+\lambda\varphi_x=0,\qquad f_y+\lambda\varphi_y=0
\] 在该点成立。因为 \(f_x\ne0\)，所以 \(\lambda\ne0\)；又
\(\varphi_y\ne0\)，故 \(f_y=-\lambda\varphi_y\ne0\)。

\begin{enumerate}
\def\labelenumi{\arabic{enumi}.}
\setcounter{enumi}{4}
\item
  证明函数
\item
  证明 \(F(x)=\int_0^\infty\frac{\sin(xt)}{1+t^2}\,dt\) 在
  \([0,\infty)\) 连续，在 \((0,\infty)\) 上具有连续导函数。
\end{enumerate}

【OCR公式整理】由 \(|\sin(xt)/(1+t^2)|\le1/(1+t^2)\)，原积分在 \(x\ge0\)
上一致收敛，故 \(F\) 连续。形式导数积分为 \[
\int_0^\infty\frac{t\cos(xt)}{1+t^2}\,dt.
\] 任取 \(\delta>0\)，\(x\ge\delta\) 时， \[
\left|\int_A^B\cos(xt)\,dt\right|\le\frac2x\le\frac2\delta.
\] 又 \(t/(1+t^2)\) 在 \(t\ge1\) 单调递减趋于零，故导数积分在
\([\delta,\infty)\) 上由 Dirichlet 判别法一致收敛。于是 \[
F'(x)=\int_0^\infty\frac{t\cos(xt)}{1+t^2}\,dt
\] 在 \([\delta,\infty)\) 连续。由 \(\delta\) 任意，结论成立。
\end{document}
